% II. Анализ на възможните решения и обосновка на избора.
% Всяко решение тук се взема по критерий, който е формулиран ПРЕДИ избора,
% за да не се окаже обосновката съчинена след факта.
\section{Анализ на възможните решения и обосновка на избора}
\label{sec:alternatives}

Изборите в този проект са четири: как се изпълнява изчислителният работник, как се
описва инфраструктурата, как се създава натоварването и как се събират метриките.
Всеки от тях се оценява по три критерия, общи за целия проект: наличие на съпоставим
еквивалент при двамата доставчици, възможност описанието да остане едно, и
пригодност на получените измервания за сравнение.

\subsection{Начин на изпълнение на изчислителния работник}

Разгледани са три възможности. Първата е \term{виртуални машини} с ръчно
инсталирано приложение, тоест чист IaaS. Тя дава пълен контрол и почти еднакво
поведение при двамата доставчици, но времето за стартиране на нов възел е от порядъка
на минути, което прави наблюдението на еластичността бавно и скъпо. Втората е
\term{управлявана услуга за контейнери} без достъп до възлите под нея. Тя намалява
административната работа, но всеки доставчик я оформя различно и еднаквото описание
става трудно. Третата е \term{Kubernetes} в управляван вид, тоест Amazon EKS и Azure
AKS, при която описанието на самото приложение е идентично и различно остава само
описанието на клъстера.

\textbf{Избран е третият вариант.} Той е единственият, при който описанието на
работното натоварване буквално съвпада при двамата доставчици, което е предпоставка
за осмислено сравнение. Цената на избора е по-висока постоянна такса за управляващия
слой и по-дълго първоначално разгръщане, и двете се отчитат при анализа на разходите.

\subsection{Инструмент за декларативно описание на инфраструктурата}

Скриптовете на командните интерфейси на самите доставчици отпадат веднага: те са
императивни, не са идемпотентни и не позволяват общо описание. Собствените средства
на доставчиците, AWS CloudFormation и Azure Bicep, са зрели, но всеки от тях работи
само при своя доставчик, което би означавало две несвързани описания. Остават
инструментите с множество доставчици: Terraform и неговото разклонение OpenTofu, и
Pulumi, при който инфраструктурата се описва на език за общо предназначение.

\textbf{Избран е Terraform, съответно OpenTofu.} Причината е, че при него разликата
между двата доставчика се локализира в отделни модули, а общата част, тоест
променливите, изходите и редът на разгръщане, остава една. Pulumi предлага същото, но
въвежда зависимост от среда за изпълнение на език, който иначе не участва в проекта.
Разликата между Terraform и OpenTofu е лицензионна, не техническа, и изборът между тях
не влияе на резултатите.

\subsection{Генератор на синтетично натоварване}

Изискването към генератора е едно: да произвежда еднакво натоварване срещу двете
разгръщания и да отчита закъснението по персентили, а не само средна стойност.
Разгледани са \code{wrk}, \code{k6} и \code{Apache JMeter}. Първият е най-лек и дава
най-малко собствено изкривяване, но описанието на сценария при него е ограничено.
Третият е най-богат, но е тежък и изисква отделна инфраструктура за самия себе си.

\textbf{Избран е k6}, защото сценарият се описва в текстов файл, който се версионира
заедно с останалата част на проекта, и защото инструментът извежда персентили
директно. \TODO{Да се потвърди, че генераторът се пуска от една и съща машина срещу
двете разгръщания, и да се опише как се отчита разликата в мрежовото разстояние до
двата региона.}

\subsection{Събиране на метрики}

Собствените средства за наблюдение на доставчиците, Amazon CloudWatch и Azure Monitor,
са налични без допълнителна работа, но измерват различни величини по различен начин и
затова не са годни за пряко сравнение. Затова метриките се събират от самото
приложение в общ формат, а средствата на доставчика се използват само за контролна
проверка.

\textbf{Избрани са Prometheus и Grafana}, разгърнати в самия клъстер. Така двете
разгръщания се наблюдават с един и същ код и една и съща конфигурация, а единствената
разлика остава средата под тях.

\subsection{Обобщение на направените избори}

\TODO{Да се добави обобщаваща таблица със структура: решение, разгледани алтернативи,
критерий, избран вариант, цена на избора. Таблицата се попълва след като изборите
бъдат потвърдени при реализацията, за да не описва намерения вместо действия.}
